% Generated draft: requires mathematical and source review.
\documentclass[11pt]{article}
\usepackage[margin=25mm]{geometry}
\usepackage{fontspec}
\setmainfont{Latin Modern Roman}
\usepackage{amsmath,amssymb,mathtools,unicode-math}
\setmathfont{Latin Modern Math}
\usepackage{enumitem,needspace,xurl,hyperref}
\setlength{\parindent}{0pt}
\setlength{\parskip}{5pt}
\setlength{\emergencystretch}{3em}
\newcommand{\N}{\mathbb{N}}
\newcommand{\Z}{\mathbb{Z}}
\newcommand{\Q}{\mathbb{Q}}
\newcommand{\R}{\mathbb{R}}
\newcommand{\C}{\mathbb{C}}
\newcommand{\F}{\mathbb{F}}
\newcommand{\A}{\mathbb{A}}
\newcommand{\PP}{\mathbb{P}}
\newcommand{\E}{\mathbb{E}}
\newcommand{\eps}{\varepsilon}
\newcommand{\dd}{\,\mathrm{d}}
\begin{document}
\section*{UnsolvedMath ID: 20002072}
\textbf{Source problem code:} catalogue entry
\subsection{Definitions and mathematical statement}
Question 1.14 is expressly terminological: ``Stein is to Weinstein as plurisubharmonic is to what?'' The source explicitly says this is not a mathematical problem. The supplied aid offers ``Lyapunov function'' for compatibility with a fixed Liouville field, ``Weinstein Morse function'' when Morse handle data are part of the structure, and sometimes ``potential'' when all surrounding data are fixed. It reports no universally imposed single word.

A classical open Weinstein structure includes a symplectic form \(\omega\), a complete Liouville vector field \(Z\), and a proper bounded-below Morse function \(\phi\) for which \(Z\) is gradient-like. The Liouville condition is \(\lambda=\iota_Z\omega\), \(\dd\lambda=\omega\). Strict Lyapunov compatibility includes \(\dd\phi(Z)>0\) away from the common critical/zero set; the chosen definition may require an additional quantitative or local gradient-like condition. Generalized-Morse variants allow prescribed birth-death singularities. For a compact domain the boundary/outward-field formulation replaces the open exhaustion/completeness conditions. A proposed name should not suppress which of these data are intended.

The aid supplies the sign-checked Stein dictionary. On a complex manifold \((W,J)\), choose \(\dd^c\phi=-\dd\phi\circ J\), assume \(\omega=\dd\dd^c\phi\) is positive, and set \(g(v,w)=\omega(v,Jw)\). Then
\[
\lambda=\dd^c\phi,\qquad \iota_{\nabla_g\phi}\omega=\lambda,
\]
so the induced Liouville vector field is the K\"ahler gradient of \(\phi\). Completeness, exhaustion and Morse properties are separate conditions, not consequences of the word ``strictly plurisubharmonic'' alone.

The aid reports three distinctions, illustrated on \((\R^2,\dd x\wedge\dd y)\). A fixed function need not determine the Liouville field: \(\phi=x^2+y^2\) is strictly increasing away from the origin along each distinct complete Liouville field
\[
Z_a=\tfrac12(x\partial_x+y\partial_y)+a(-y\partial_x+x\partial_y),\qquad a\in\R.
\]
An exhausting Morse function need not be Lyapunov for a specified field: for \(Z_0\), the function \((x-1)^2+y^2\) has \(\dd\phi(Z_0)<0\) at \((1/2,0)\). Conversely an exhausting strict Lyapunov function need not be Morse or generalized Morse: take \(\phi=F(x^2+y^2)\), with \(F(s)=\int_0^s e^{-1/u}\,\dd u\) and the integrand extended by zero at \(u=0\). Its critical point is infinitely flat, while \(\dd\phi(Z_0)>0\) elsewhere. These are distinctions among properties, not a new nomenclature theorem.

The trailing sentence introducing Takao Akahori belongs to the next question and supplies no additional task here.
\subsection{Short English statement}
The source asks for terminology for the function in a Weinstein structure and explicitly disclaims an open mathematical problem. The supplied aid distinguishes Lyapunov compatibility, Morse data, exhaustion and Liouville completeness, with a sign-fixed plurisubharmonic dictionary.
\subsection{Sources}
\begin{itemize}
\item Question 1.14, 2003 AIM workshop \emph{Holomorphic curves in contact geometry}: \url{https://aimath.org/WWN/contactgeom2/contactgeom2.pdf}. The source itself calls the item a question of terminology; the Takao Akahori transition begins the next item.
\item Local registry \texttt{data/unsolvedmath/problems.json}, permanent ID 20002072, code \texttt{catalogue entry}; canonical \texttt{aim-geometry-notes.json}, \texttt{notes[409]}; updated 14 August 2026. The unrecoverable statement-status label is retained despite its verification reproducing the terminological question. The original literature field is empty.
\item Attempt 1: context only, problem status at run not a problem, rigor label proved, novelty confidence low; current registry open/open triage. The aid's terminology discussion has no identified full primary-source entries. It is machine generated, not peer reviewed, and requires independent verification before citation.
\end{itemize}
\end{document}
